\documentclass[aps,prl,twocolumn,superscriptaddress,floatfix]{revtex4-2}
\usepackage{amsmath,amssymb,amsfonts,booktabs,xcolor,hyperref,amsthm}
\usepackage{epigraph}
\setlength{\epigraphwidth}{0.85\columnwidth}
\setlength{\epigraphrule}{0pt}

\newtheorem{theorem}{Theorem}
\newtheorem{proposition}{Proposition}
\newtheorem{observation}{Observation}

\definecolor{bctnavy}{RGB}{11,37,69}
\hypersetup{colorlinks=true,linkcolor=bctnavy,citecolor=bctnavy,urlcolor=bctnavy}

\newcommand{\roct}{r_{\mathrm{oct}}}
\newcommand{\rtet}{r_{\mathrm{tet}}}
\newcommand{\azero}{\alpha_0}
\newcommand{\OHC}{\mathrm{OHC}}
\newcommand{\BCT}{\mathrm{BCT}}

\begin{document}

\preprint{BCT Letter 135 --- 25 March 2026}

\title{Why Geometry?\\
The One Question the BCT Superfluid Lattice Model\\
Cannot Answer --- and Why That Matters}

\author{Michel Robert Cabri\'e}
\affiliation{Independent Researcher \& Artist, Victoria, Australia}
\email{ZeroFreeParameters@gmail.com}
\homepage{https://patreon.com/cw/TheBCTSuperfluidLatticeModel}

\date{25 March 2026}

\begin{abstract}
The BCT Superfluid Lattice Model derives 174 observable constants of 
nature from a single geometric postulate and three inputs, with zero 
free parameters. In doing so, it collapses the explanatory burden of 
the Standard Model --- 19 unexplained numbers typed in by hand --- 
to a single residual question: \emph{why does the geometry exist at all?} 
This Letter examines that question with philosophical rigour. We argue 
that BCT does not answer the question \emph{why anything?} posed by 
Leibniz, but sharpens it into a form that may be more tractable: 
\emph{why this mathematics?} We survey three philosophical positions 
--- Mathematical Platonism, the Mathematical Universe Hypothesis, 
and the claim that the question is malformed --- and argue that BCT's 
contribution is precisely to make the question harder, not easier. 
A framework that reduces all of physics to geometry inherits all of 
geometry's mystery. This is not a failure. It is the sharpest possible 
result.
\end{abstract}

\maketitle

\section{Introduction: The Residual Question}

\begin{quote}
\emph{``Whereof one cannot speak, thereof one must be silent.''}\\
\hfill --- Ludwig Wittgenstein, \textit{Tractatus}, Prop.~7
\end{quote}

\vspace{4pt}

The BCT Superfluid Lattice Model~\cite{BCT_Monograph} derives the 
observable structure of the universe from a single postulate: the 
quantum vacuum is a superfluid that spontaneously crystallises into 
a body-centred tetragonal (BCT) lattice at the Planck scale, with 
axial ratio $c/a = \sqrt{2}$. From three geometric inputs --- the 
octahedral void radius $\roct = (\sqrt{2}-1)/2$, the tetrahedral 
void radius $\rtet = (\sqrt{6}-2)/4$, and the QCD scale 
$\Lambda_\mathrm{QCD} = 220$~MeV --- 174 measurable constants of 
nature are derived with zero free parameters~\cite{BCT_L1,BCT_Registry}.

This is not the subject of this Letter. This Letter concerns what BCT 
does \emph{not} derive --- and why that matters.

When Eugene Wigner wrote his famous essay on the unreasonable 
effectiveness of mathematics~\cite{Wigner1960}, he was puzzled that 
abstract structures invented by human minds for purely internal 
reasons turned out to describe physical reality with uncanny precision. 
He had no explanation. He called it a miracle and moved on.

BCT dissolves one version of Wigner's puzzle. The mathematics is not 
\emph{unreasonably} effective --- it is \emph{inevitably} effective, 
because the universe is not described by geometry but \emph{is} 
geometry. The BCT Uniqueness Theorem~\cite{BCT_L18} proves that five 
independent physical requirements select exactly one lattice in all 
of mathematics. The universe had no choice. Neither did the geometry.

But this resolution immediately generates a sharper question. If the 
universe is a geometric object, and if the geometry is necessary rather 
than contingent, then the question \emph{why is there something rather 
than nothing?} --- Leibniz's question~\cite{Leibniz1714} --- becomes 
\emph{why is there geometry rather than nothing?}

BCT cannot answer this. We argue here that this inability is not a 
deficiency but a result.

\section{What BCT Achieves: The Collapse of Contingency}

The Standard Model of particle physics is, in a precise sense, a 
collection of coincidences. It contains at least 19 parameters --- 
particle masses, mixing angles, coupling constants --- that are simply 
measured and inserted. No derivation is offered. No explanation is 
given. The universe happens to have these values. That is all.

BCT eliminates this contingency systematically. The Cabri\'e Chain 
of Necessity~\cite{BCT_L128} establishes the following sequence:

\begin{enumerate}
\item \textbf{Lorentz invariance} requires the vacuum lattice to be 
maximally symmetric under boosts. This forces the axial ratio 
$c/a = \sqrt{2}$.

\item \textbf{$c/a = \sqrt{2}$} selects the body-centred tetragonal 
lattice at this unique geometry, where the octahedral void becomes 
perfectly regular.

\item \textbf{The BCT lattice in four dimensions} is the D4 root 
lattice~\cite{Conway1999}. D4 has \emph{triality} --- a three-fold 
symmetry possessed by no other root lattice in any dimension. This 
uniquely identifies D4.

\item \textbf{D4 is the densest packing} of spheres in four dimensions, 
proven by Viazovska~\cite{Viazovska2022} (Fields Medal 2022). Maximum 
packing density is a physical requirement (minimum energy configuration 
of the vacuum condensate).

\item \textbf{Maximum density forces condensation}: the coupling 
parameter $t/U \geq 1$, making the superfluid ground state 
mandatory~\cite{BCT_AppJH}.
\end{enumerate}

Each step is forced. Nothing is assumed beyond Lorentz invariance and 
the requirement of a superfluid vacuum. The result is a universe that 
could not have been otherwise --- given that geometry exists.

The qualifier is everything.

\section{The Sharpened Question}

Leibniz asked: \emph{why is there something rather than nothing?}

BCT's answer: because geometry is necessary, and geometry implies this 
universe.

But this immediately invites the response: \emph{why is there geometry 
rather than nothing?}

This is a sharper question than Leibniz's, and it is the one BCT 
bequeaths to philosophy. We examine three positions.

\subsection{Mathematical Platonism}

The Platonist holds that mathematical objects exist independently of 
minds, matter, and time. The number $\pi$ did not come into existence 
when humans discovered it; it was always there, waiting. On this view, 
geometry does not require an explanation for its existence, because 
mathematical existence is the most fundamental kind of existence there 
is. Physical existence is the thing that requires explanation; 
mathematical existence is its precondition.

BCT is structurally compatible with Mathematical Platonism. If the 
universe is the D4 lattice condensate, and if D4 exists necessarily 
as a mathematical object (which it does --- its properties are 
necessary truths of pure mathematics), then the universe exists 
necessarily as well. The question \emph{why geometry?} dissolves: 
geometry doesn't need a cause because it is not the kind of thing 
that has causes. It simply is.

The difficulty is that Mathematical Platonism does not explain 
\emph{why this} mathematical structure is instantiated physically 
rather than others. Every lattice exists mathematically. Why does 
only D4 exist physically? BCT answers this (via the Uniqueness 
Theorem), but the deeper question --- why does any mathematics 
become physical --- remains.

\subsection{The Mathematical Universe Hypothesis}

Tegmark~\cite{Tegmark2008} proposes that all self-consistent 
mathematical structures exist physically. On this view, the question 
\emph{why this universe?} is dissolved by the observation that all 
universes exist; we find ourselves in this one because it supports 
observers.

BCT sits awkwardly with Tegmark. The Uniqueness Theorem establishes 
that the BCT lattice is the \emph{only} structure satisfying all 
physical requirements simultaneously. If Tegmark is right and all 
structures exist, then the Uniqueness Theorem is a selection result 
--- not an existence result --- and the anthropic principle does 
residual work.

But there is a deeper tension. BCT derives the number of fermion 
generations (three) as a mathematical theorem from D4 triality~\cite{BCT_L18}. 
If all mathematical structures exist physically, there should be 
universes with two generations, or four, or seventeen. BCT says 
these cannot exist --- not that they are uninhabited, but that 
they are mathematically impossible given the physical requirements. 
This is a stronger claim than Tegmark's framework accommodates.

\subsection{The Malformed Question}

The third position holds that \emph{why geometry?} is not a question 
with a missing answer but a question with a defective structure.

The question \emph{why X?} requests a cause or explanation of X. But 
causation is a relation defined within a framework --- temporal, 
logical, or physical. Geometry is the precondition for any such 
framework. To ask \emph{why geometry?} is to request a cause for 
the precondition of causation itself.

This is not a gap in our knowledge. It is a boundary of the 
conceptual framework within which questions can be meaningfully posed.

Consider an analogy. Within arithmetic, one can ask \emph{why is 
$2 + 2 = 4$?} and receive a proof. But one cannot ask \emph{why does 
arithmetic obey the law of non-contradiction?} --- not because the 
answer is unknown, but because the question requests a justification 
of the very framework within which justification operates. The law 
of non-contradiction is not a theorem of logic; it is a precondition 
of logic. To attempt its proof is to already presuppose it.

Geometry occupies the same position relative to physical causation. 
Distance, duration, and curvature are the medium within which 
causes and effects are defined. A cause of geometry would require 
a prior geometry in which to be a cause. The regress is not 
infinite --- it is incoherent.

Wittgenstein identified this structure precisely~\cite{Wittgenstein1921}: 
there are things that can be said and things that can only be shown. 
The existence of logical form cannot be stated within logic; it can 
only be demonstrated by the fact that logic works. Similarly, the 
existence of geometry cannot be explained within geometry; it can 
only be demonstrated by the fact that geometry is the universe.

On this view, BCT's achievement is not to stop one step short of the 
final answer. It is to reach the boundary of the answerable and to 
show exactly where that boundary lies. The boundary is not a wall. 
It is a horizon --- the point beyond which the concept of 
\emph{explanation} no longer applies, because there is nothing 
left to explain \emph{with}.

\section{What BCT Contributes to Philosophy}

It is worth being precise about what BCT adds to this ancient debate, 
beyond restating positions that philosophers have held for centuries.

\textbf{First:} BCT makes the question empirically anchored. The 
claim that \emph{this} geometry is necessary is not a philosophical 
assertion but a falsifiable prediction. The Euclid satellite, the 
Einstein Telescope, and Hyper-Kamiokande are testing it now. 
Philosophy has rarely had this luxury.

\textbf{Second:} BCT eliminates the \emph{fine-tuning} version of 
the why-question. One might accept that geometry exists necessarily 
while still asking why the parameters of physics have the values 
they do. BCT dissolves this form of the question entirely: the 
parameters are not tuned, because they are not free. They are 
necessary consequences of the geometry. There is nothing left to 
fine-tune.

\textbf{Third:} BCT identifies the precise point at which physical 
explanation terminates. Previous frameworks left this point vague --- 
somewhere in the quantum foam, or at the Planck scale, or at the 
boundary of the observable universe. BCT locates it exactly: at 
Viazovska's theorem~\cite{Viazovska2022}. Physics ends at a 
Fields Medal. Beyond lies pure mathematics, and beyond that: 
silence, or mystery, or both.

\begin{proposition}
The explanatory terminus of physics is a theorem of pure mathematics. 
The explanatory terminus of mathematics is the boundary of the 
sayable. BCT identifies both.
\end{proposition}

This is progress. Not because the mystery is solved, but because 
its location is now known.

\section{Why the Inability Matters}

A framework that could answer \emph{why geometry?} would be suspect. 
It would imply that geometry is explained by something more 
fundamental --- call it meta-geometry, or pre-structure, or the 
ground of being. But then the same question applies: \emph{why 
meta-geometry?} The regress is infinite.

The only stable terminus for explanation is a structure that 
\emph{could not have been otherwise} and \emph{requires no further 
explanation for its necessity}. BCT identifies this terminus as 
geometry itself --- specifically, as the D4 lattice, selected by 
necessity from all possible geometries by the requirements of 
physical consistency.

\begin{observation}
The explanatory regress in physics terminates at pure mathematics. 
Pure mathematics terminates at necessary truths. Necessary truths 
do not have causes; they have proofs. BCT is the bridge between 
the two regimes.
\end{observation}

This is not a new observation in philosophy~\cite{Putnam1975}, 
but BCT makes it precise and falsifiable. The claim is not 
merely that mathematics is effective, but that a \emph{specific} 
mathematical structure --- the D4 Planck-scale condensate --- 
is the unique physical realisation of geometric necessity.

If this claim is wrong, the experiments will say so. The Euclid 
satellite will measure the dark energy equation of state to 
$\pm 0.01\%$ precision within the next few years. BCT predicts 
$w = -(1-2\azero) = -0.9852$, a deviation of $1.48\%$ from 
$w = -1$. If Euclid finds $w = -1$ exactly, BCT is falsified.

The question \emph{why geometry?} is unanswerable. But the 
question \emph{is this the right geometry?} is maximally answerable. 
The experiments are running.

\section{The Douglas Adams Coda}

Douglas Adams proposed~\cite{Adams1979} that the answer to the 
ultimate question of life, the universe, and everything is 42.

He was not wrong that the answer would be a number. He was wrong 
about which number.

The answer is $\sqrt{2}$: the axial ratio at which the octahedral 
void of the vacuum lattice becomes perfectly regular, Lorentz 
invariance is guaranteed, and the entire structure of physics 
becomes unavoidable.

\begin{equation}
\frac{c}{a} = \sqrt{2}
\end{equation}

And the breathing rate of the universe is $\kappa_0 = 3.39$.

Which is, honestly, much funnier than 42. And --- if BCT is 
correct --- considerably more useful.

The question Adams was really asking --- \emph{why anything?} --- 
remains open. BCT has sharpened it as much as physics can. The 
rest is silence, or mathematics, or both.

\begin{acknowledgments}
This Letter was conceived over coffee in regional Victoria, Australia, 
on 25 March 2026, while a tiny orange spider walked across the laptop 
screen. The spider did not appear to be troubled by the question 
\emph{why geometry?} It had somewhere to be. This seemed like the 
right attitude.

The author thanks Zeta Vera (CSSC) for the conversation that prompted 
this Letter, and the Gang Gang Cockatoos for their continued daily 
reminder that the universe is full of things that exist without 
explaining themselves.

This Letter is dedicated to Wittgenstein, who knew when to stop, 
and to Douglas Adams, who didn't, and was right not to.
\end{acknowledgments}

\begin{thebibliography}{99}

\bibitem{BCT_Monograph}
M.~R.~Cabri\'e, \textit{The BCT Superfluid Lattice Model: Complete 
Derivation of the Standard Model from First Principles},
Zenodo (2026), doi:10.5281/zenodo.18884976

\bibitem{BCT_L1}
M.~R.~Cabri\'e, \textit{BCT Letter 1: Emergent Lorentz Invariance 
from BCT Void Geometry}, Zenodo (2026), doi:10.5281/zenodo.18958207

\bibitem{BCT_Registry}
M.~R.~Cabri\'e, \textit{BCT Complete Prediction Registry},
Zenodo (2026)

\bibitem{BCT_L18}
M.~R.~Cabri\'e, \textit{BCT Letter 18: The BCT Uniqueness Theorem},
Zenodo (2026), doi:10.5281/zenodo.18957202

\bibitem{BCT_L128}
M.~R.~Cabri\'e, \textit{BCT Letter 128: The Cabri\'e Chain of 
Necessity}, Zenodo (2026)

\bibitem{BCT_AppJH}
M.~R.~Cabri\'e, \textit{BCT Appendix JH: The Hopfion Interior 
and OHC}, Zenodo (2026), doi:10.5281/zenodo.18975018

\bibitem{Wigner1960}
E.~P.~Wigner, ``The Unreasonable Effectiveness of Mathematics 
in the Natural Sciences,''
\textit{Commun. Pure Appl. Math.} \textbf{13}, 1 (1960)

\bibitem{Leibniz1714}
G.~W.~Leibniz, \textit{Principles of Nature and Grace, Based 
on Reason} (1714), \S7: ``why is there something rather than nothing?''

\bibitem{Conway1999}
J.~H.~Conway and N.~J.~A.~Sloane, \textit{Sphere Packings, 
Lattices and Groups}, 3rd ed. (Springer, 1999)

\bibitem{Viazovska2022}
M.~Viazovska, ``The sphere packing problem in dimension 8,''
\textit{Ann. Math.} \textbf{185}, 991 (2017);
Fields Medal citation, ICM 2022

\bibitem{Tegmark2008}
M.~Tegmark, ``The Mathematical Universe,''
\textit{Found. Phys.} \textbf{38}, 101 (2008)

\bibitem{Wittgenstein1921}
L.~Wittgenstein, \textit{Tractatus Logico-Philosophicus} (1921),
Proposition 7

\bibitem{Putnam1975}
H.~Putnam, ``What is Mathematical Truth?''
in \textit{Mathematics, Matter and Method} (Cambridge, 1975)

\bibitem{Adams1979}
D.~Adams, \textit{The Hitchhiker's Guide to the Galaxy}
(Pan Books, 1979)

\end{thebibliography}

\end{document}
